\label{sec:conclusion}

This paper has documented open-seat creation and retirement incentives under \textsc{bisect}
algorithmic redistricting and compared these outcomes to enacted 2022 congressional maps
and to available neutral comparators. The central findings are consistent across NC, WI,
and TX: algorithmic maps produce approximately twice as many open seats as enacted maps,
and they produce retirement incentive scores that are substantially higher for all three
states.

The comparison to commission-drawn maps---where available---suggests that algorithmic
open-seat rates are comparable to neutral human redistricting outcomes, consistent with
the general proposition that algorithmic redistricting produces incumbency consequences
similar to what a neutral human process would produce, not more disruptive.

From a legal and policy perspective, the open-seat analysis highlights a fundamental
tradeoff in redistricting design. Protecting incumbents through engineered home districts
reduces electoral competition, increases retirement security for sitting members, and
produces delegations that are more stable but less responsive to constituent preferences
and partisan tides. Algorithmic redistricting, by being indifferent to these protections,
produces more open seats, more competitive elections, and more responsive delegations---
as a geometric side effect, not as a policy goal.

Courts evaluating algorithmic redistricting as a remedy or standard for future maps
should understand that the algorithmic map's open-seat count is not an excessive
disruption to incumbency---it is the disruption that geographic chance would produce
in the absence of deliberate protection. The enacted map's lower open-seat count
represents the premium that legislators have extracted from the redistricting process
to insulate themselves from competitive elections. Whether that premium is constitutionally
permissible is a question for courts; this paper establishes the empirical magnitude
of the premium.
